\RequirePackage{etex}

% Vi benytter 2 sidet, a4 størrelse og dansk sprog. 
\documentclass[11pt,a4paper,oneside,danish]{memoir}
\usepackage[paper=A4,pagesize]{typearea}
 \usepackage[utf8]{inputenc} % danske tegn

% Choose language by uncommenting danish or english
\usepackage[danish]{babel} % dansk sprogpakke
% \usepackage[english]{babel}


%% Pakker til fonts, farver, referencer og ligenende er placeret relevant sted i dokumentet.
     
     %% FONT
     \usepackage[default,osfigures,scale=0.95]{opensans} % Alternatively use the option 'defaultsans' instead of 'default' to replace the sans serif font only.
     \usepackage[T1]{fontenc}
     
     % Common packages 
     \usepackage{footnote} % Package used to controle footnote(specially in tables etc.)
     \usepackage{amsmath}
     \usepackage{slantsc}
     \usepackage{multirow} % Merge table rows
     \usepackage{longtable} % Table that can span multiple pages
     \usepackage{geometry}
     \usepackage[section]{placeins} % \FloatBarrier
     \usepackage[compatibility=false]{caption}
     \captionsetup{font=small,labelfont=bf,textfont=it}
     \usepackage{graphicx}
     \usepackage{afterpage}
     \usepackage[draft]{fixme} %hvor 'options' bl.a. inkluderer final, draft, danish, inline og margin
     \usepackage{rotating}
     \usepackage{eso-pic} % til inds�ttelse af forside.
     \usepackage{rotating} % rotate stuff(EVERYTHING!11)

	% % MATH STUFF 
    %\usepackage{algorithmic}
    %\usepackage{harpoon} %Tilf�jer harpun-pile over st�rrelser i math, m.m.
    %\usepackage{pst-plot} % For axes   
    %\usepackage{breqn}
    %\usepackage{fourier}  %................... Roman+math - Utopia
    
  
    %\usepackage{nomentbl} % http://ctan.mackichan.com/macros/latex/contrib/nomentbl/nomentbl.pdf

    %\usepackage{psfrag} % insæt tag, mest brugt i forhold til grafer

%%% DIVERSE CUSTOMS PAKKET NED I FILER FOR OVERBLIK
   % \input{styles/HaxTilPreamble.tex}
   % \input{styles/CustomFunctions.tex}
%	\input{Styles/code.tex} 
    \usepackage{todonotes}


% new commands for individually colored todo 
\newcommand{\todosander}[2]{
\todo[#1, color=green!40]{#2 \\- Sander}
}
\newcommand{\todosimon}[2]{
\todo[#1, color=red!40]{#2 \\- Simon}
}
\newcommand{\todorasmus}[2]{
\todo[#1, color=yellow!40]{#2 \\- Rasmus}
}
\newcommand{\todokristoffer}[2]{
\todo[#1, color=purple!40]{#2 \\- Simon}
}


% Color package. Use with listings package and use for color a table
%\usepackage[usenames,dvipsnames,svgnames,table]{xcolor}

    %%%%%%%%%%%%%%%%%%%%%%%%%%%%%%%%%%%%%%%%%%%%%%%%%%%%%%%%%%%%%%%%%%%%%%%%%%%%%
%                            Custom for funktioner                          %
%                     			      	        		                    %
%%%%%%%%%%%%%%%%%%%%%%%%%%%%%%%%%%%%%%%%%%%%%%%%%%%%%%%%%%%%%%%%%%%%%%%%%%%%%


% Package for wrapping figures in text
\usepackage{wrapfig}

% Set the path where LaTeX looks for pictures.
\graphicspath{{figure/}}

% You need a newsubfloat element to use subcaption
\newsubfloat{figure}

% Command to set caption styles
\captionnamefont{\bfseries\small}
\captiontitlefont{\itshape\small}
\subcaptionlabelfont{\bfseries\small}
\subcaptionfont{\itshape\small}

% When we don't use chapter numbering, 
% the figure numbering are shown like 
% fx 0.3 if we don't use this line.
\counterwithout{figure}{chapter}


% Command \figur{filename}{caption}{label}{width} 
% for inserting a new figure.
\newcommand{\figur}[4]{ 
 \begin{figure}[ht] 
 \centering 
 \includegraphics[width=#4\textwidth]{#1} 
 \caption{#2} 
 \label{#3} 
 \end{figure}
 } 
 
% Command \wrapfigur{filename}{caption}{label}{width}{r/l}
% for inserting a new figure wraped in the text. 
\newcommand{\wrapfigur}[5]{ 
 \begin{wrapfigure}{#5}{#4 \textwidth}
   \begin{center}
    \includegraphics[width=#4\textwidth]{#1}
   \end{center}
   \caption{#2}
   \label{#3} 
 \end{wrapfigure}
}
 
% Command \subfigur{filename1}{filename2}{caption}
% {subcaption1}{subcaption2}{label1}{label2}
% {width1}{width1} for inserting a new figure.
\newcommand{\subfigur}[9]{
\begin{figure}[ht]
\centering
\subbottom[{#4}\label{#6}]%
    {\includegraphics[width=#8\textwidth]{#1}}\hfill
\subbottom[{#5}\label{#7}]%
    {\includegraphics[width=#9\textwidth]{#2}}
\caption{#3}
\label{fig:subfigures}
\end{figure}
}
% Command \truefigure{filename}{caption}{label}
% for inserting a picture without scaling.
\newcommand{\truefigure}[3]{ 
 \begin{figure}[ht] 
 \centering 
 \includegraphics{#1} 
 \caption{#2} 
 \label{#3} 
 \end{figure}
 } 

% Keeps floats in the section. 
% Help control place to put pix
\usepackage[section]{placeins}

% Commands for rightaligning wraped text in a tabular
\newcommand{\rr}{\raggedright}
\newcommand{\tn}{\tabularnewline}

\usepackage{multirow}

%Tabular pakke som Wex har fundet, men som vi ikke bruger alligvel
%\usepackage{tabu}

% Include code pieces from files
% Documentation:
% ftp://ftp.tex.ac.uk/tex-archive/macros/latex/contrib/listings/listings.pdf
\usepackage{listings}


%Package with enumirate command with extra options (used within cells of tabular)
\usepackage[inline]{enumitem}
    
    % add a box using the themecolor. Can be used to wrap around a chapintro for introducing an assignment
    %use \begin{assignment} and end\{assignment} to box it
    %http://texdoc.net/texmf-dist/doc/latex/mdframed/mdframed.pdf
     
    \usepackage[xcolor]{mdframed} 
    \newmdenv[linecolor=ThemeColor]{assignment}
%    \usepackage{Placeins}
    
% % % % % % % % % % % % % % ********% % %Generelt design % % % % % % % % % % % % % % % % % % %
 
%\usepackage[table,dvipsnames]{xcolor} % 'table' er i stedet for \usepackage{colortbl} og 'dvipsnames' er et farveskema hvor man angiver farver med navne

\usepackage{xcolor} % farver kan indsætttes i henhold til RGB med definecolor

%Valg af temafarve % Denne farve bruges i chapterstyle %
\definecolor{ThemeColor}{RGB}{209, 67, 26}

% Valg af font %
\usepackage[sc]{mathpazo}
\linespread{1.05}         % Palatino needs more leading (space between lines)
\usepackage[T1]{fontenc}

% Pakker inkludere overordnede designvalg. Pakkerne er valgt for overskuelighed %
	 %%%%%%%%%%%%%%%%%%%%%%%%%%%%%%%%%%%%%%%%%%%%%%%%%%%%%%%%%%%%%%%%%%%%%%%%%%%%%
%                             Chapter and Section                           %
%                     			      Style	       		                    %
%%%%%%%%%%%%%%%%%%%%%%%%%%%%%%%%%%%%%%%%%%%%%%%%%%%%%%%%%%%%%%%%%%%%%%%%%%%%%


% % % % Redefined so that header will show even on chapter start. % % % %
\makeatletter
    \let\stdchapter\chapter
    \renewcommand*\chapter{%
    \@ifstar{\starchapter}{\@dblarg\nostarchapter}}
    \newcommand*\starchapter[1]{%
        \stdchapter*{#1}
        \thispagestyle{fancy}
        \markboth{\MakeUppercase{#1}}{}
    }
    \def\nostarchapter[#1]#2{%
        \stdchapter[{#1}]{#2}
        \thispagestyle{fancy}
    }
\makeatother




\usepackage{titlesec} % Pakke til opsætning af chapter og sektioner


% % defining new font for the chapternumber. Calling it FuckingHuge 
\newcommand{\FuckingHuge}{% Needs a 'stretchable' font
      \usefont{\encodingdefault}{\rmdefault}{b}{n}%
      \fontsize{35}{45}\normalfont\color{ThemeColor}} %60 80


%%% TAL I OVERSKRIFTER, både chapters og sections. 
    \maxsecnumdepth{subsubsection} %subsubsection er dybdeniveau for numering.
   	\setsecnumdepth{subsubsection}


  
 % % % % CHAPTER OPSÆTNING % % % % % % % % % % % % % % % % % % % % %

%Hang thoose numbers! I margin med tallene for både sections og chapters! En lækker chapter style!
\chapterstyle{hangnum}
 
%Modificering af ovenstående style
%Placering af chapter overskriften på siden i forhold til sidens start %
\renewcommand*{\chapterheadstart}{\vspace*{-20pt}}


% FARVE PÅ CHAPTER % %
%TItel farve og font valg
\renewcommand*{\chaptitlefont}{\normalfont\Huge\color{ThemeColor}}
%Numering farve og font valg
\renewcommand{\chapnumfont}{\normalfont\FuckingHuge\color{ThemeColor}}



% % % % SECTION OPSÆTNING % % % % % % % % % % % % % % % % % % % %
\hangsecnum % vi smidder numering ud i marken

%Vi benytter tema farven til numeringen
\renewcommand\thesection{\color{ThemeColor}\thechapter.\arabic{section}}



    
    


	 \input{Styles/header.tex}
	 
% Setup the margins:
\setlrmarginsandblock{1in}{*}{1.5}
\setulmarginsandblock{3cm}{4cm}{*}

% Setup header and footer: %
% \setheadfoot{headheight}{footskip}
% \setheaderspaces{headdrop}{headsep}{ratio}
\setheadfoot{45pt}{1cm}
\setheaderspaces{1cm}{*}{*}
\checkandfixthelayout
% Se dette pdf for dokmentaton:
% www.tug.dk/_media/foredrag/2003-12-03/slides-src.pdf

 \checkandfixthelayout[nearest]
    \checkthelayout
    \fixthelayout
    \baselineskip \topskip

%BibTex setup
\usepackage[semicolon,authoryear,square]{natbib}

\bibliographystyle{plainnat}
%%%%%%%%%%%%%%%%%%%%%%%%%%%%%%%%%%%%%%%%%%%%%%%%%%%%%%%%%%%%%%%%%%%%%%%%%%%%%%%%%%%%%%%%%%%%%%

% % Smart listing function. Documentation: http://texblog.org/tag/makenomenclature/
\usepackage{nomencl}
\makenomenclature 

% % % % % % % % % % % % % % % % Reference opsætning % % % % % % % % % % % % % % % % % % % % % % %

\usepackage{varioref} % for references, http://www.ctex.org/documents/packages/bibref/varioref.pdf
   
% External reference package, zref makes the references into hyperlinks.
\usepackage{url}
\usepackage{hyperref}
%\usepackage{zref-xr}
\usepackage{xr}

% % % % % % % % % % % % % % % % % % % % % % % % % % % % % % % % % % % % % % % % % % % % % % % % % % %

%\usepackage{showlabels}

% footruleskip is allready defined in memoir, so we have to undefine it, in order to include fancyhdr for our footer and header.
\let\footruleskip\undefined
\usepackage{fancyhdr}% For making fancy headers http://ctan.org/pkg/fancyhdr

\AtBeginDocument{\renewcommand\contentsname{Indholdsfortegnelse}}

% % % % % % % % % % % % % Inkluderinger af ønskede kapitler % % % % % % % % % % % % %

\begin{document}
\newgeometry{left=3cm,bottom=0.1cm}

\centerline{\Huge\huge\bfseries\color{ThemeColor}Inversion of Control} 
\vspace{1em}
\centerline{\large\color{ThemeColor}Aarhus School of Engineering - Aarhus Universitet}


\centerline{\large\date{10. April 2015}}

%\color{Black}
%\vspace{8em} % Tilføj nogle tomme linjer
\begin{center}
    \begin{tabular}{ l l p{6cm} }
	 \textbf{Gruppe 11}& & \\
	& & \\
    Stud. nr.: 201270115 & Johnny Davidsen Poulsen & \\\hline
    & & \\
    Stud. nr.: 201271039  & Kristian Rossen Kristensen & \\\hline
     \end{tabular}
      	
\end{center}
\thispagestyle{empty} % No header and footer on this page
\restoregeometry
\setcounter{tocdepth}{2}
\tableofcontents

 

% Setup header and footer
\pagestyle{fancy}
\fancyhf{} % Clear all header an footer fields 
\renewcommand{\headrulewidth}{0pt} % Delete headruler
\fancyhead[R]{\includegraphics[height=40pt]{header}} % Insert picture in left side of header
%\fancyhead[R]{\rightmark} % insert chapter title in right side of header
\fancyfoot[LE,RO]{\thepage} % insert pagenumber on right side of even page, left side of odd page 
%\setcounter{page}{1}

\thispagestyle{fancy}
\chapter{Beskrivelse}
\label{chp:description}

\section{Indhold}
Inversion of Control er et design pattern som er vidt udbredt i eventdrevne systemer hvor klassernes afhængigheder af hinanden ikke nødvendigvis er fastlagt compile time. 

IoC design patterenet implementeres som oftests gennem et framework. Frameworket søger for at danne et stykke kode, en såkaldt IoC container, som får tildelt flow control run time. Blandt de mest udbredte frameworks til .NET kan nævnes Autofac, Unity, Castle og Ninject \footnote{Scott Hanselman \href{http://www.hanselman.com/blog/ListOfNETDependencyInjectionContainersIOC.aspx}{Link}}. IoC frameworks anvendes også i vid udstrækning på Android\footnote{leftshift \href{http://leftshift.io/android-inversion-of-control-dependency-injection-dagger-part-1}{Link}}, iOS\footnote{stackoverflow \href{http://stackoverflow.com/questions/8782398/recommended-ioc-framework-for-ios}{Link}} og andre platforme som anvender event driven programming.

\figur{koncept}{Konceptet bag IoC}{fig:concept}{0.5}

De fleste frameworks har det til fælles at de anvender dependency injection. En IoC container konfigureres i kode og delegeres flow control i en bootstrapper eller .config fil eller lignende når programmet startes. Herefter er det containerens ansvar at instantiere klasser runtime efterhånden som programmet har behov for dem. 

\section{Formål}
IoC anvender dependency injection til at skabe mindre afhængighed mellem generiske klasse i et program. Inversion of control tager kontrollen fra den enkelte klasse og skaber et design pattern som er mere hensigtsmæssig en dependency inversion og factory patterens i systemer som er eventdrevne og hvor alle afhængigheder i et system derfor ikke er bestemte compile time. IoC implementeres som en container i et stykke af koden som overlades flow control.

\section{Konsekvenser}
Konsekvensen ved at bruge inversion of control er uafhængige klasser som er fuldstændigt afkoblede. Klasserne har ingen indbyrdes afhængihed og fungerer derfor nærmere som er service andre klasser kan anvende. IoC gør ens program mere skalerbar. Det gør klasserne nemmere at genbruge i andre programmer. Klassernes indbyrdes uafhængighed gør desuden programmet mere egnet i et event drevet system i modsætning til sekventielt afviklet kode.

\thispagestyle{fancy}
\chapter{Sammenligning}
\label{chp:comparison}

\section{Problemet uden IoC}

For at forstå IoC design princippet, er et godt udgangspunkt problematikkerne der kan opstå, hvis man ikke benytter sig af dette. Nedenstående eksempel består af en Customer klasse som indeholder et Address objekt. Problemet er den tætte kobling mellem Customer og Address - hvis Address klassen på noget tidspunkt ændrer sig, vil dette også have indflydelse på kompileringen af Customer klassen.

\figur{Customer1.png}{Customer og Address problemet}{fig:CustomerOgAddress}{0.4}

Det primære problem, som skaber den tætte kobling er, at Customer klassen skal oprette et Address objekt.
Endvidere er Customer klassen koblet direkte op mod Address klassen, så tilføjelse af eks. HomeAddress eller OfficeAddress vil ikke kunne lade sig gøre.

Generelt ses det også, at hvis Customer ikke kan oprette et Address objekt vil det ikke være muligt, at konstruere et Customer objekt.


\section{Problemets generelle løsning}

Problemerne opstår generelt når oprettelsen af et objekt foregår inde i en klassen. Det vil derfor, i Customer Address eksemplet, være ideelt at flytte oprettelsen af Address objektet fra Customer klassen til en anden klasse. I stedet for at Customer har kontrollen af oprettelse af objekter, gives denne til en anden klassen. Dette er den primære ide bag design princippet Inversion of Control (IoC).
Dette princip er også defineret som Hollywood princippet: "Don't call us we will call you" eller fra Address objektets synsvinkel: "Don't create me I will create my self using some one else".


\section{Specifikke måder at implementere IoC}

IOC er et design princip, som i sig selv ikke løser problemerne. Det er her Dependency Injection (DI) kommer ind i billedet, som er en måde at implementere IOC på. 

\figur{IOC.png}{IOC i relation til DI}{fig:IOCandDI}{0.6}
\FloatBarrier

Der er fire måder hvor på DI kan implementeres: 

\begin{itemize}
	\item Constructor
	\item Getter og Setter
	\item Interface
	\item Service locator
\end{itemize}

Alt afhængig af hvordan DI benyttes i programmet, kan de forskellige DI metoder benyttes. Neden for ses et ekempel med Constructor Dependency Injection. I vores tilfælde vil Address blive oprettet af klienten og blot sendt med som constructor parameter til Customer. Nu har Customer en reference til Address og kan frit benytte sig af denne.

\figur{DI.png}{Customer og Address constructor DI}{fig:constructorDI}{0.5}

\section{Factory}

Den faktiske måde DI bliver implementeret på er også forskellig. Det er muligt, at oprette alle de afhængige objekter i en enkelte klasse, eksempelvis main(), og sende dem videre ned i de oprettede objekter. Det primære problem er, at når oprettelsen af objekter sker i en enkelt klasse opstår der alligevel høj kobling. Eksempelvis, når en concrete klasse skal udvides med en ekstra constructor parameter, skal alle overstående klasser implementere denne. Dette kan dog løses ved hjælp af Factory design mønstret. 

\figur{Factory.png}{Factory design pattern}{fig:constructorDI}{0.5}

Her sendes der blot et Factory objekt rundt, som kan opfylde oprettelsen af de enkelte objekters behov. Selv om Factory kan løse mange af problemerne, er der dog stadig argumenter for, at dette ikke er et optimalt design. Det største problem med en Factory er, at den kan ikke genbruges i andre applikationer. Alle valg om oprettelse er med henblik på den nuværende applikation, og det er sjældent muligt, at klasser portere til en anden applikation. Endvidere er det et stort arbejde for compileren, at gennemføre runtime udskiftning af de enkelte Factories.


\section{Container}

En Container kan tage ansvaret for at oprette, konfigurere og holde styr på objekterne. Det er muligt, at konfigurere objekterne i Containeren, som også hjælper med at oprette og injecte objekterne i de respektive klasser. 

I tilfældet med Customer og Address, vil Containeren være den, som registrer de to objekter separat og injecter Address objekt ind i Customer klassen. Disse Containers er ofte tilbudt som et framework, der gør det let at implementere og har et højt abstraktionsniveau. I afsnit 3 fremstilles eksemplificeres implementeringen af en IoC container.
\thispagestyle{fancy}
\chapter{Eksempler}
\label{chp:examples}
I dette afsnit vises eksempler på anvendelsen af to forskellige frameworks, et til .NET og et til Android. 
\section{Autofac}
Autofac er et af det mest anvendte frameworks til .NET. Autofac tilbyder IoC containers som via Dependency Injection, injecter alle afhængigheder til programmets klasser via deres constructorere. Autofac kan installeres som NuGet packages eller ved manuelt at tilføje Autofac's library filer til projektet.

\figur{examples/AutofacNugetInstall}{Installation af Autofac via NuGet}{fig:installautonuget}{0.7}

\subsection{Konfiguration}
Når Autofac er tilføjet til projektet skal containeren tildeles flow control over programmet. Et konsolprogram tilføjer dette ganske som forventet ved at configurere en container i starten af Main() som vist i figur \ref{fig:configureiocconsole}. Herefter køres metoden WriteData som er en metode man selv implementerer. WriteData sætter containers scope til applikationens scope.

\figur{examples/AutofacConsoleAppRegister}{Konfiguration af IoC Container i konsol applikation}{fig:configureiocconsole}{0.5}
\FloatBarrier

I WPF har man ikke adgang til programmets Main(). Her anvendes istedet programmets onStartup metode i Application klassen til at tildele flow control til IoC containeren. Et eksempel på dette ses i figur \ref{fig:configureiocwpf}

\figur{examples/AutofacWPFRegister}{Konfiguration af IoC Container i WPF applikation}{fig:configureiocwpf}{0.5}

Ved konfigurationen af sin container registrer man alle de klasser som containeren skal kunne injecte. Metoden er typisk at angive hvilken konkret klasse som containeren skal implementere for et interface. På den måde kan man hurtigt og enkelt skifte mellem implementeringer som det kendes fra et abstract factory pattern.

\subsection{Injection}
Efter at man har tildelt flow control til sin IoC container og registreret alle sine klasser som skal injectes, begynder den egentlige magi ved Autofac. Alle klasser instantieres automatisk af containeren, når der er behov for dem. Autofac sørger desuden for at instantiere og constructor injecte alle parametre som en klasse har.

Dette betyder, at Autofac er fuldstændig usynlig og arbejder i baggrunden når man skriver sine klasser. Man søger bare for at angive alle sine afhængiheder som constructorparametre som Autofac herefter forsyner klassen med ved instantiering.

\section{RoboGuice}
RoboGuice er et Denpendency Injection framework til Android, som også kan konfigureres til at anvende IoC containers. RoboGuice tilføjes til et projekt ved at importere referencerne i ens build.gradle eller pom.xml fil hvis man bruger Maven. Et eksempel på import i gradle kan ses i figur \ref{fig:importroboguice}

\figur{examples/RoboGuiceImport}{Import af RoboGuice i Android projekt}{fig:importroboguice}{0.45}

\FloatBarrier
\subsection{Konfiguration}
Da Android applikationer er eventdrevne, ligner opsætningen meget den som anvendes med Autofac i WPF. Man starter med at lave en custom implementering af Application klassen. Heri registrer man sin container og knytter denne til en Modul klasse hvori man kan konfigurere containeren og binde klasser med interfaces.

I figur \ref{fig:robocontainerreg} ses det hvorledes at Application classens onCreate klasse overrides og containeren tildeles flow control. Klassen RoboModule registreres til konfiguration af containeren.

\figur{examples/RoboGuiceRegisterContainer}{Registrering af contrainer i Application}{fig:robocontainerreg}{0.6}

Implementeringen af RoBoModule ses i figur \ref{fig:roboclassreg} hvor klassen Foo bindes til interfaceet IFoo. Hver gang IFoo injectes i en klasse vil containeren således forsyne klassen med en reference til et Foo objekt.

\figur{examples/RoboGuiceRegisterClasses}{Registrering af klasser i container}{fig:roboclassreg}{0.9}

\FloatBarrier
\subsection{Injection}
Syntaxen for injection anvender Android Annotation (AA) til at property injecte afhængighederne i en klasse. I figur \ref{fig:roboinject} ses det hvorledes at klassen Foo injectes med syntaxen \emph{@Inject}. Bemærk at det er interfacet IFoo som er angivet da dette er bundet til klassen Foo gennem konfigurationen i RoboModule som vist i forrige afsnit. På den måde ville det således være enkelt at skifte alle afhængiheder af Foo ud med eksempelvist Moo, ved blot at ændre bindingen.

\figur{examples/RoboGuiceInject}{Injection af klasser med RoboGuice}{fig:roboinject}{0.55}

\thispagestyle{fancy}
\chapter{Konklusion}
\label{chp:Konklusion}
\todo[inline]{Konkludere paa emnet – hvornår giver det bedst mening at anvende emnet, og i hvilke(n) situation(er) giver det ingen mening?}

%\bibliographystyle{plainnat}
%\bibliography{bibliografi}% Selects .bib file AND prints bibliography
\thispagestyle{fancy}

%\listoftodos % Make a list of todo's
\end{document}